Fix an identifying cell, a compatible \(P_{\mathrm{obs}}\) on that cell, \(\mathcal M\in\mathfrak F(\mathcal G,P_{\mathrm{obs}})\), and \((a,y)\). By the forward direction of \(\text{thm:response-fiber-equivalence}\), there is \(\theta_{\mathcal M}\) satisfying
\[
\mathsf C_{\mathcal G}(\theta_{\mathcal M};P_{\mathrm{obs}}),
\qquad
q_{a,y}(\theta_{\mathcal M})=Q_a^{\mathcal M}(y).
\tag{8}
\]
On an identifying cell, \(\mathsf I_{a,y}\) is true. Its defining negated existential says that any two compatible response vectors have equal \(q_{a,y}\). Therefore \(\mathsf V_{a,y}(P_{\mathrm{obs}},z)\) has a unique value \(z_*\). By \(\text{def:tcf-procedure}\), if no rational trace validates on the whole cell, the stored functional is this unique-value graph. If a rational trace \(N/D\) is stored, its whole-cell verification gives \(D\neq0\) and \(Dz=N\) at \(z=z_*\), so the rational annotation also equals \(z_*\); the division is licensed by the verified nonzero denominator. Hence in both cases (8) gives
\[
f_{a,j}(P_{\mathrm{obs}})(y)=z_*=q_{a,y}(\theta_{\mathcal M})=Q_a^{\mathcal M}(y).
\]
This argument is semantic over the real numbers and uses no numerical encoding of the coordinates of \(P_{\mathrm{obs}}\).

Apply the equality first with \(a=a_1\), then with \(a=a_0\), multiply the two equalities by \(s_Y(y)\), and sum over the finite alphabet \(\mathcal X_Y\). The definition of \(\tau_{s_Y}\) gives
\[
\sum_y s_Y(y)\{f_{a_1,j}(P_{\mathrm{obs}})(y)-f_{a_0,j}(P_{\mathrm{obs}})(y)\}
=\tau_{s_Y}.
\]